\begin{figure*}[t]
\centering
\scriptsize
\setlength{\fboxsep}{1em}
\setlength{\fboxrule}{1pt}
\fcolorbox{black}{white}{
\parbox{0.955\textwidth}{
\begin{alltt}


You will be provided instructions for completing a task followed by a task to complete.

Instructions:

1. Understand the Dataset: The IMDb Movie Reviews dataset contains 50,000 reviews from the Internet Movie Database (IMDb). These reviews are labeled as either positive or negative, with an even number of each. The sentiment of the reviews is determined based on the score given by the reviewer. A review is labeled as negative if it has a score of 4 or less out of 10, and it is labeled as positive if it has a score of 7 or more out of 10. The dataset does not include more than 30 reviews per movie.

2. Analyze the Passage: Read the passage carefully. This will be a review from the IMDb dataset. Pay attention to the language used, the tone of the review, and any explicit or implicit expressions of sentiment towards the movie. 

3. Determine the Sentiment: Based on your analysis of the review, determine whether the sentiment is positive or negative. Remember that positive reviews generally contain praise for the movie, its plot, characters, cinematography, etc., and have a score of 7 or more out of 10. Negative reviews, on the other hand, generally contain criticism of the movie and have a score of 4 or less out of 10.

4. Answer the Question: Based on your determination of the sentiment, select the appropriate output: 'Negative' or 'Positive'.

Remember, the goal is to correctly identify the sentiment of the review based on the text provided.

Passage: Lemuel Gulliver (Ted Danson) is a doctor who goes missing at sea, leaving pregnant wife Mary (Mary Steenburgen) behind. Eight years later, he turns up, disheveled and seemingly mad - babbling about his adventures in the lands of the tiny Lilliputians, the giant Brobdingnags, the floating island of the intellectual Laputa, and the Houyhnhnms, a race of intelligent, talking horses who have to deal with the Yahoos - a race of bestial humans - among many other adventures. The not-so-good Dr. Bates (James Fox), who has designs on Lemuel's wife, has Gulliver incarcerated in a mental institution, and Lemuel, Mary, and son Thomas (Tom Sturridge) must find a way to prove his sanity.<br /><br />A splendid adaptation of Jonathan Swift's satirical novel, this film is a magnificent adaptation on so many levels: the story, the satire, the characters, the visuals, the brilliant cast. It's simply a treat to watch, and it's almost amazing considering that it was a made-for-TV film.<br /><br />The film does a brilliant job of capturing Swift's vicious satire, which cuts like a hatchet through British society of the time, but still resonates today. The wise Brobdingnags and the Houyhnhnms are almost perfect individuals who find it virtually impossible to understand why Gulliver speaks with such pride of the vices and corruptions of his society. The scenes where Gulliver struggles to prove himself different from the Yahoos are perhaps the best, with biting satire in describing how they pick their leaders ("they seem to pick the worst among them. . . who rules until they find someone even worse"), go to war ("We only go to war for a very good reason - such as they are weaker than us, or we want all of their land"), etc. The scenes involving Laputa are also effectively done - the intellectuals are so wrapped up in their specialized fields that they have no time for anything else, and really possess little common sense. And the addition of the asylum plot line enhances the story greatly - Dr. Bates is truly nasty character, and when he gives a speech to the inquiry on Gulliver's alleged vices, it's quite clear that he's describing his own faults.<br /><br />The film makes use of beautiful, and fairly convincing CGI effects depicting the very diverse settings of the novel with great effect. The contrast of sizes is done in a very skillful way, and all of the worlds depicted in the story are convincing in their own way. The cinematography (particularly that concerning the asylum) and the costumes are brilliantly done. The editing of the present with Lemuel's memories is a device which could be awkward, but works very well.<br /><br />The cast is truly wonderful; a veritable who's-who of British and American talent. Ted Danson gives an excellent, multi-layered performance as Gulliver, showing effectively his transformation from a person bewildered by his strange surroundings, to the lunatic state he was in when he reappears, to his rational, intellectual personality at the end. Most well-known for his work on sit-com, Danson shows that he's more than just Sam Malone with this wonderful serio-comic performance. Mary Steenburgen is effective as his wife, and James Fox is absolutely repulsive as Bates. The rest of the cast is made up mostly of cameos, with Peter O'Toole, Omar Sharif, Warwick Davis, Kristin Scott Thomas, Geraldine Chaplin, Alfre Woodward, Edward Fox, and Sir John Gielgud being the most memorable - but even the smallest parts are very well-played.<br /><br />While not 100%

Follow the instructions to generate an explanation that reasons towards the correct answer to the task above. End the explanation with the correct answer. [/INST] Explanation: 

The review provided is a positive review of the movie "Gulliver's Travels." The reviewer uses positive language to describe the film, such as "magnificent adaptation," "brilliant cast," and "triumph of story and images." They also highlight the film's ability to capture the satire of Jonathan Swift's novel, which is a key aspect of the movie. Additionally, the reviewer praises the film's visuals, costumes, and cinematography.

The only potential negative comment in the review is a brief mention of the film's deviation from the book, but the reviewer quickly dismisses this by stating that it is "not 100%

Therefore, based on the language, tone, and content of the review, it is clear that the sentiment of the review is positive.

Answer: Positive

Therefore, the correct label among {'Negative', 'Positive'} (just the label) to the original task is below.

Sentiment: Positive
\end{alltt}
}}
\caption{A complete running example of IMDB.}
\label{fig:imdb_example}
\end{figure*}
